\section{Observations: Information Revealed Along the Memory Path}
\label{sec:observations}

We use the same 14 workloads as the evaluation to examine when a request
reveals information that can remove or combine memory-system work. O1
establishes the end-to-end cost. O2 and O3 motivate local paired HBM access,
O4 and O5 motivate requester-side remote aggregation, and O6 motivates
selective reuse in the existing requesting-GPM L2.

\noindent\textbf{O1: Cost of a Full Memory Path.}
We measure each component from request arrival until it forwards the request
or response, charge local communication to the receiving component, and group
RDMA processing and wafer traversal as Remote Communication. We count only
completed demand-read MSHR leaders and exclude L1 and L2 MSHR waiting. As
Fig.~\ref{fig:obs-full-path} shows, an HBM-reaching local path takes
74--563~ns, while the seven workloads with a qualifying remote demand path
take 0.45--7.19~$\mu$s. The local L2 demand-read miss rate ranges from 2.8\%
to 100\%, and the dominant latency component changes across workloads. The
bottleneck is therefore an end-to-end path property rather than the property
of one fixed cache, network, or memory component.

\begin{figure}[!t]
  \centering
  \includegraphics[width=\columnwidth]{../../akkalat/results/2026-07-22-baseline-complete-observation/observation-analysis/figures/o1_memory_path_latency.pdf}
  \caption{Normalized full-path latency and local-L2 demand-read miss rate for
  local (L) and remote (R) MSHR leaders. MSHR waiting is excluded, and
  ``--'' denotes no qualifying remote demand path.}
  \label{fig:obs-full-path}
\end{figure}

\noindent\textbf{O2: Local Misses Expose Adjacent Demand Pairs.}
We record each real L2 read miss and locate the nearest earlier miss to the
adjacent cacheline in the same aligned access region. Within 16 L2 cycles,
13 of 14 workloads expose such a relation for 5.0\%--50.0\% of their misses.
Weighted across all workloads, adjacent relations cover 7.3\%, 15.0\%,
22.0\%, and 33.0\% of misses within 0, 4, 16, and 64 cycles, respectively
(Fig.~\ref{fig:obs-adjacent-window}). KMeans exposes no relation in the
measured window, showing why adjacent access should be learned from real
demands rather than generated unconditionally.

\begin{figure}[!t]
  \centering
  \includegraphics[width=\columnwidth]{../../akkalat/results/2026-07-22-baseline-complete-observation/observation-analysis/figures/o2_adjacent_line_cdf.pdf}
  \caption{Fraction of real L2 read misses whose adjacent cacheline miss
  arrived within a short preceding window.}
  \label{fig:obs-adjacent-window}
\end{figure}

\noindent\textbf{O3: Physical Mapping Makes Adjacent Work HBM-Friendly.}
We apply the controller's bank, row, column, and aligned-access mapping to
physical read arrivals. Within 16 controller cycles, 29.1\% of reads find an
earlier read in the same aligned HBM access unit, while only 6.2\% find a
different column in the same row (Fig.~\ref{fig:obs-dram-locality}). Thus,
the strongest local relation is not arbitrary spatial proximity but the
specific adjacent-line relation already aligned with the controller's access
unit. M1 learns this relation from completed demand pairs and attempts a
paired access only after exact cache and pending state remove an already
available half. Otherwise, the original singleton request proceeds.

\begin{figure}[!t]
  \centering
  \includegraphics[width=\columnwidth]{../../akkalat/results/2026-07-22-baseline-complete-observation/observation-analysis/figures/o3_dram_physical_locality.pdf}
  \caption{CDF of physical-read proximity for an aligned HBM access unit,
  same-row accesses, row changes within a bank, and different-bank accesses.}
  \label{fig:obs-dram-locality}
\end{figure}

\noindent\textbf{O4: Remote Requests Amplify Work Before Owner Coalescing.}
Across ten workloads, 713,568 sampled remote operations traverse requester
RDMA, the wafer network, and owner-GPM RDMA before any owner-L2 MSHR can merge
their memory work. One logical byte produces 6.05 network byte-hops, and mean
remote-transaction latency ranges from 0.465 to 7.787~$\mu$s
(Fig.~\ref{fig:obs-remote-amplification}). An owner-L2 MSHR may still reduce
later HBM accesses, but it acts only after requester endpoint and forward
network work has already been paid. Removing redundancy at requesting-GPM
RDMA therefore eliminates work that owner-side coalescing cannot recover.

\begin{figure}[!t]
  \centering
  \includegraphics[width=\columnwidth]{../../akkalat/results/2026-07-22-baseline-complete-observation/observation-analysis/figures/o4_remote_amplification.pdf}
  \caption{Network byte-hops per logical remote byte and mean end-to-end
  remote-operation latency. ``--'' denotes no remote operation in the O4
  trace window.}
  \label{fig:obs-remote-amplification}
\end{figure}

\noindent\textbf{O5: Requester RDMA Reveals Two Aggregation Scopes.}
Among 558,671 sampled remote reads, 144,474 (25.9\%) have an exact same-line
request already in flight. After removing those duplicates, 129,074 of the
remaining 414,197 reads (31.2\%) find a different line with the same
requester, owner GPM, and 4-KB page within 16 RDMA cycles
(Fig.~\ref{fig:obs-remote-aggregation}). Exact coalescing and same-page packet
grouping together cover 49.0\% of the original reads, but their contributions
vary sharply by workload. The interval characterizes available concurrency.
M2 joins only an existing batch and never delays a request for a future
partner.

\begin{figure}[!t]
  \centering
  \includegraphics[width=\columnwidth]{../../akkalat/results/2026-07-22-baseline-complete-observation/observation-analysis/figures/o5_remote_aggregation.pdf}
  \caption{Exact in-flight same-line redundancy and same-page,
  different-line aggregation opportunities at requesting-GPM RDMA.}
  \label{fig:obs-remote-aggregation}
\end{figure}

\noindent\textbf{O6: Remote Reuse Is Selective and Existing L2 Has Headroom.}
Only 17.7\% of unique remote lines recur, yet reads after the first touch
account for 33.7\% of all remote reads. The second access establishes reuse
and can trigger an L2 fill, so third-or-later accesses, which account for
22.0\% of remote reads, are the direct demand-only opportunity to avoid
another wafer traversal. The existing requesting-GPM L2 meanwhile has
60.2\% mean free capacity, with workload means ranging from 9.8\% to 94.1\%
(Fig.~\ref{fig:obs-remote-reuse}). M3 therefore retains only returns supported
by recurrence or multiple-waiter evidence, while exact replacement checks
protect local data and no new cache level is required.

\begin{figure}[!t]
  \centering
  \includegraphics[width=\columnwidth]{../../akkalat/results/2026-07-22-baseline-complete-observation/observation-analysis/figures/o6_remote_reuse_l2_headroom.pdf}
  \caption{Recurring remote lines, reads after first touch, and mean free
  capacity in the existing requesting-GPM L2. ``--'' denotes no remote read.}
  \label{fig:obs-remote-reuse}
\end{figure}

\noindent\textbf{Synthesis: Visibility Evolves Along the Request Lifetime.}
The observations reveal a sequence rather than three unrelated
optimizations. Real local misses expose temporal adjacency, physical mapping
reveals HBM compatibility, requesting-GPM RDMA exposes cross-L1 duplication
and common destinations, and completed remote activity reveals recurrence.
No single point observes all of this information. CuPath therefore acts at
the first boundary where each relation becomes actionable, removes redundant
work before it travels farther, and immediately preserves the baseline path
when no opportunity exists.
\FloatBarrier
